\documentclass[11pt]{article}
\usepackage{preambulo}
\addbibresource{bibliografia.bib}

\begin{document}

\portada{Modelos de aprendizaje automático}

\tableofcontents
\newpage

%============================================================================
\section{Planteamiento del problema de aprendizaje}
%============================================================================

El presente documento describe los modelos de aprendizaje automático aplicados al conjunto de
datos SUPPORT2 \parencite{support2}, cuyo preprocesamiento se documentó previamente. El objetivo
consiste en predecir la variable \texttt{hospdead}, que indica si un paciente fallece durante la
hospitalización analizada.

La variable objetivo admite únicamente dos valores y se encuentra registrada para la totalidad de
los casos. El problema corresponde por tanto a una tarea de \textbf{clasificación binaria
supervisada}: a cada paciente debe asignársele una de dos categorías, y el aprendizaje se realiza
a partir de ejemplos cuya categoría es conocida.

El conjunto empleado está conformado por 7{,}284 registros de entrenamiento y 1{,}821 de prueba,
con 40 variables de entrada y la variable objetivo. Todas las variables son numéricas y no
contienen valores ausentes, condición establecida durante el preprocesamiento.

Una característica del conjunto condiciona la interpretación de los resultados: las clases están
desbalanceadas. Los fallecimientos representan el 25.9\% de los registros, de modo que un
clasificador que asignara siempre la categoría mayoritaria acertaría en el 74.1\% de los casos sin
utilizar información alguna.

%============================================================================
\section{Criterio de selección de los modelos}
%============================================================================

Se evaluaron tres clasificadores en lugar de uno solo. La razón es que el desempeño de un modelo
aislado no permite distinguir entre dos situaciones distintas: que el algoritmo sea adecuado para
el problema, o que la señal predictiva esté presente en los datos con independencia del algoritmo.
Comparar modelos que deciden mediante mecanismos diferentes permite separar ambos efectos.

Los tres seleccionados responden a mecanismos de decisión distintos entre sí:

\begin{itemize}
\item La \textbf{regresión logística} construye una frontera lineal: pondera cada variable y
      combina las ponderaciones en una única puntuación.
\item El \textbf{árbol de decisión} construye reglas: formula preguntas encadenadas sobre una
      variable por vez.
\item El método de \textbf{vecinos más cercanos} no construye regla alguna: clasifica cada caso
      según los registros que más se le parecen.
\end{itemize}

Se incorporó además un cuarto modelo de referencia, que no aprende de los datos y predice siempre
la clase mayoritaria. Su función se detalla en la Sección \ref{sec:referencia}.

La Tabla \ref{tab:configuracion} resume la configuración de los cuatro.

\begin{table}[H]
\centering
\caption{Configuración de los modelos evaluados}
\label{tab:configuracion}
\small
\begin{tabular}{p{3.6cm} p{5.4cm} p{1.6cm}}
\toprule
\textbf{Modelo} & \textbf{Parámetros} & \textbf{Escalado} \\
\midrule
Regresión logística & Penalización L2, hasta 3000 iteraciones & Sí \\
Árbol de decisión & Profundidad máxima 5, criterio Gini & No \\
Vecinos más cercanos & 25 vecinos, distancia euclidiana & Sí \\
Referencia trivial & Predice siempre la clase mayoritaria & No \\
\bottomrule
\end{tabular}
\end{table}


%============================================================================
\section{Regresión logística}
%============================================================================

\subsection{Mecanismo de decisión}

La regresión logística asigna un coeficiente a cada variable de entrada, calcula la suma ponderada
de los valores del paciente y transforma dicha suma en una probabilidad mediante la función
logística de la Ecuación \eqref{eq:logistica}:

\begin{equation}
P(\text{fallece}) = \frac{1}{1 + e^{-(\beta_0 + \beta_1 x_1 + \cdots + \beta_n x_n)}}
\label{eq:logistica}
\end{equation}

El resultado se interpreta directamente como probabilidad de fallecimiento. Un coeficiente
positivo indica que valores mayores de esa variable incrementan el riesgo estimado; uno negativo,
que lo reducen.

\subsection{Configuración}

Se empleó la penalización L2, que limita la magnitud de los coeficientes y previene que el modelo
asigne un peso desproporcionado a una variable concreta. El número máximo de iteraciones se fijó
en 3000, con holgura suficiente: sobre estos datos la optimización converge en 17 iteraciones.

La estandarización de las variables resulta necesaria en este caso, y el motivo es precisamente la
penalización. Al castigar los coeficientes de mayor magnitud, una variable medida en unidades
pequeñas requeriría un coeficiente elevado para ejercer el mismo efecto, y la penalización la
afectaría de manera desproporcionada respecto de las demás.

\subsection{Variables de mayor peso}

La Figura \ref{fig:coeficientes} presenta los coeficientes de mayor magnitud obtenidos sobre las
variables estandarizadas. La estandarización permite compararlos entre sí, dado que todos operan
sobre la misma escala.

\begin{figure}[H]
\centering
\includegraphics[width=0.85\textwidth]{coeficientes.png}
\caption{Coeficientes de mayor magnitud de la regresión logística. En rojo los que incrementan el
riesgo estimado, en azul los que lo reducen.}
\label{fig:coeficientes}
\end{figure}

Los dos coeficientes de mayor peso corresponden a \texttt{sps} y \texttt{aps}, los puntajes de
fisiología medidos al tercer día, seguidos por el indicador de coma y el puntaje \texttt{scoma}.
El modelo concede por tanto el mayor peso a las mediciones de gravedad clínica, resultado
coherente con el fenómeno que se pretende predecir.

Entre los coeficientes negativos figuran varios indicadores de grupo diagnóstico, en particular
insuficiencia cardíaca congestiva, cáncer de colon y enfermedad pulmonar obstructiva crónica. Su
signo indica que, a igualdad de las restantes variables, pertenecer a esos grupos reduce el riesgo
estimado de fallecimiento durante la hospitalización respecto de la categoría de referencia.

%============================================================================
\section{Árbol de decisión}
%============================================================================

\subsection{Mecanismo de decisión}

El árbol de decisión clasifica mediante preguntas encadenadas sobre una variable por vez. En cada
nodo selecciona la variable y el umbral que mejor separan las dos clases, evaluados mediante el
índice de Gini, y repite el procedimiento sobre cada una de las dos ramas resultantes.

A diferencia de la regresión logística, el resultado puede representarse gráficamente y leerse
como una secuencia de reglas, lo que constituye su principal ventaja en contextos donde la
decisión debe justificarse.

\subsection{Profundidad}

La profundidad máxima se fijó en cinco niveles. Un árbol sin restricción de profundidad continúa
dividiendo hasta aislar casos individuales, con lo que reproduce el conjunto de entrenamiento en
lugar de generalizar. El límite de cinco niveles produce un árbol de 32 hojas, cantidad que
conserva la posibilidad de inspección.

\subsection{Estructura obtenida}

La Figura \ref{fig:arbol} presenta los dos primeros niveles del árbol entrenado.

\begin{figure}[H]
\centering
\includegraphics[width=\textwidth]{arbol.png}
\caption{Primeros dos niveles del árbol de decisión. La profundidad completa es de cinco niveles
y 32 hojas.}
\label{fig:arbol}
\end{figure}

La primera división utiliza \texttt{aps} con un umbral de 50.5, y las divisiones subsiguientes
recurren a \texttt{scoma} y \texttt{sps}. La regla inicial puede enunciarse directamente: los
pacientes cuyo puntaje de fisiología APACHE III supera 50.5 presentan un riesgo sustancialmente
mayor.

\subsection{Ausencia de estandarización}

El árbol se entrenó sobre las variables en sus unidades originales, sin estandarizar. La
estandarización resta la media y divide por la desviación, operación que no altera el orden
relativo de los valores. Dado que el árbol únicamente compara cada valor contra un umbral,
selecciona la misma división y produce idénticos resultados sobre datos estandarizados o sin
estandarizar.

Aplicar la transformación no modificaría el resultado y sugeriría una dependencia que no existe.
Se omite por ese motivo.

%============================================================================
\section{Vecinos más cercanos}
\label{sec:vecinos}
%============================================================================

\subsection{Mecanismo de decisión}

El método de vecinos más cercanos no construye un modelo durante el entrenamiento: almacena los
registros disponibles. Para clasificar un paciente nuevo, localiza los $k$ registros más próximos
según la distancia euclidiana y asigna la categoría mayoritaria entre ellos.

\subsection{Necesidad de estandarización}

En este método la estandarización no constituye una mejora sino un requisito. Como muestra la
Ecuación \eqref{eq:distancia}, la distancia euclidiana agrega las diferencias de todas las
variables en un único valor:

\begin{equation}
d(a, b) = \sqrt{\sum_{i=1}^{n} (a_i - b_i)^2}
\label{eq:distancia}
\end{equation}

Sin estandarización, las variables de mayor recorrido dominan la suma. La glucosa abarca un
intervalo de aproximadamente 318 unidades, mientras que el pH se desplaza entre 7.25 y 7.60: una
diferencia de 50 unidades de glucosa pesaría miles de veces más que la diferencia máxima
observable en pH. La proximidad quedaría determinada casi exclusivamente por una variable.

\subsection{Selección del número de vecinos}

El valor de $k$ se determinó mediante validación cruzada sobre el conjunto de entrenamiento,
evaluando varios valores y comparando el AUC promedio. La Figura \ref{fig:vecinos} presenta el
resultado.

\begin{figure}[H]
\centering
\includegraphics[width=0.78\textwidth]{vecinos.png}
\caption{AUC promedio en validación cruzada según el número de vecinos.}
\label{fig:vecinos}
\end{figure}

El procedimiento se realizó sobre el conjunto de entrenamiento y no sobre el de prueba. Determinar
un parámetro a partir del conjunto reservado para evaluación invalidaría la evaluación posterior,
dado que el valor seleccionado habría sido elegido precisamente por su desempeño sobre esos datos.

\subsection{Comportamiento observado}

El AUC aumenta de forma sostenida con el número de vecinos: 0.762 con cinco vecinos, 0.817 con
veinticinco y 0.833 con doscientos uno. La ausencia de un máximo dentro del intervalo evaluado
merece interpretación.

El comportamiento es atribuible a la dimensionalidad del conjunto. Con 40 variables de entrada,
las distancias entre registros tienden a homogeneizarse y la noción de proximidad pierde capacidad
discriminante. El modelo compensa esa pérdida promediando sobre vecindarios cada vez más amplios,
procedimiento que se aproxima progresivamente a estimar la proporción global de la clase.

Se adoptó $k = 25$ por constituir un valor intermedio que conserva el carácter local del método.
La elección implica renunciar a aproximadamente 0.016 puntos de AUC respecto del mejor valor
evaluado.

%============================================================================
\section{Modelo de referencia}
\label{sec:referencia}
%============================================================================

Se incorporó un cuarto modelo que no utiliza las variables de entrada y predice invariablemente la
clase mayoritaria. Su propósito no es competir con los anteriores sino establecer el umbral por
debajo del cual ningún resultado puede considerarse informativo.

Dado el desbalance del conjunto, este modelo alcanza una exactitud de 0.741. Cualquier
clasificador que no supere ese valor no aporta información, con independencia de que su exactitud
pueda parecer elevada en términos absolutos. Su inclusión evita interpretar como logro un
desempeño que se obtiene sin aprender nada.

%============================================================================
\section{Protocolo de entrenamiento}
%============================================================================

El entrenamiento se realizó exclusivamente sobre el conjunto de entrenamiento, reservando el de
prueba para la evaluación documentada en el informe correspondiente.

La comparación entre modelos se efectuó mediante validación cruzada estratificada de diez
particiones. El procedimiento divide el conjunto de entrenamiento en diez bloques que conservan la
proporción de clases, entrena con nueve y evalúa sobre el restante, repitiendo la operación diez
veces. El promedio resultante es más estable que el obtenido de una partición única. La selección
del número de vecinos, documentada en la Sección \ref{sec:vecinos}, empleó cinco particiones por
tratarse de una búsqueda sobre un único modelo.

El escalado se integró como etapa inicial de un flujo conjunto con el modelo, y no como
transformación previa aplicada a los archivos. La distinción no es formal: dentro de la validación
cruzada, el escalador debe ajustarse sobre los cuatro bloques de entrenamiento de cada iteración y
aplicarse al bloque restante. Si se ajustara una sola vez sobre el conjunto completo, cada
evaluación se realizaría con un escalador que ya incorporó información del bloque evaluado.

La semilla de los procedimientos aleatorios se fijó en 42, de modo que los resultados son
reproducibles.

La Tabla \ref{tab:resultados} resume el desempeño obtenido sobre el conjunto de prueba. Su
análisis detallado, junto con las matrices de confusión y las pruebas de significancia
estadística, se desarrolla en el informe de evaluación.

\begin{table}[H]
\centering
\caption{Desempeño de los modelos sobre el conjunto de prueba}
\label{tab:resultados}
\small
\begin{tabular}{p{3.8cm} r r r r}
\toprule
\textbf{Modelo} & \textbf{AUC} & \textbf{Exactitud} & \textbf{Precisión} & \textbf{Recall} \\
\midrule
Regresión logística & 0.879 & 0.827 & 0.736 & 0.519 \\
Árbol de decisión & 0.847 & 0.807 & 0.701 & 0.447 \\
Vecinos más cercanos & 0.835 & 0.797 & 0.721 & 0.356 \\
Referencia trivial & 0.500 & 0.741 & 0.000 & 0.000 \\
\bottomrule
\end{tabular}
\end{table}


%============================================================================
\section{Convergencia de los tres modelos}
%============================================================================

Los tres clasificadores emplean mecanismos de decisión distintos y, pese a ello, se apoyan en las
mismas variables.

La regresión logística concede sus mayores coeficientes a \texttt{sps} y \texttt{aps}. El árbol de
decisión selecciona \texttt{aps} para su primera división y \texttt{scoma} y \texttt{sps} para las
siguientes. La selección univariada de variables, documentada en el informe de preprocesamiento,
sitúa a esos mismos tres puntajes entre los más informativos.

La coincidencia resulta relevante porque los tres procedimientos operan bajo criterios
independientes: uno optimiza coeficientes de una combinación lineal, otro maximiza la separación
de clases en divisiones sucesivas, y el tercero evalúa la asociación individual de cada variable
con el objetivo. Que converjan en los mismos puntajes de severidad indica que la señal predictiva
reside en los datos y no en la elección del algoritmo.

El resultado es además coherente desde la perspectiva clínica: \texttt{aps}, \texttt{sps} y
\texttt{scoma} son precisamente los índices construidos para cuantificar la gravedad del paciente.

%============================================================================
\section{Conclusiones}
%============================================================================

Se entrenaron tres clasificadores con mecanismos de decisión diferentes sobre el conjunto
preprocesado, acompañados de un modelo de referencia que establece el umbral mínimo de
desempeño informativo.

La regresión logística obtuvo el mayor AUC, con 0.879. El árbol de decisión alcanzó 0.847 y el
método de vecinos más cercanos 0.835. Los tres superan la referencia trivial, cuyo AUC es 0.500 y
cuya exactitud de 0.741 ilustra por qué esta última métrica resulta insuficiente en conjuntos
desbalanceados.

Cada modelo presenta un perfil distinto. La regresión logística combina el mejor desempeño con la
posibilidad de examinar el peso asignado a cada variable. El árbol admite una lectura directa como
secuencia de reglas, a costa de aproximadamente tres centésimas de AUC. El método de vecinos
resulta el menos apropiado para este conjunto, no por configuración deficiente sino porque la
dimensionalidad degrada la noción de proximidad en que se sustenta.

La convergencia de los tres en los mismos puntajes de severidad constituye el hallazgo de mayor
relevancia: indica que el desempeño obtenido proviene de la información contenida en los datos y
no de la elección del procedimiento.

La determinación de si las diferencias observadas entre los tres modelos son estadísticamente
significativas requiere pruebas específicas, que se desarrollan en el informe de evaluación.

\printbibliography[title={Referencias}]

\end{document}
